\noindent\markforreview
\vspace{1em}

A circle has center \(O\), and points \(A\) and \(B\) lie on the circle.
The measure of arc \(AB\) is \(60^\circ\), and the length of this arc is 4 inches.
What is the circumference, in inches, of the circle?

\vspace{0.75em}

\choicebox{A}{4}
\choicebox{B}{8}
\choicebox{C}{16}
\choicebox{D}{24}

\vspace{2em}
\newpage

\noindent\markforreview
\vspace{1em}

In the \(xy\)-plane, the graph of function \(f\), where \(y = f(x)\), has exactly 8 \(x\)-intercepts.
One of these \(x\)-intercepts is \((17, 0)\).
The rational function \(g\) is defined by:
\[
g(x) = \frac{f(x)}{x - 17}
\]
In the \(xy\)-plane, how many \(x\)-intercepts does the graph of \(y = g(x)\) have?

\vspace{0.75em}

\choicebox{A}{5}
\choicebox{B}{6}
\choicebox{C}{7}
\choicebox{D}{8}

\vspace{2em}
\newpage

\noindent\markforreview
\vspace{1em}

A researcher investigated two species of mites: a predator and its prey.
At the start of a week, there was an equal number of the two species.
At the end of the week:

\begin{itemize}
\item The number of prey had increased by 2700\% of the number of prey at the start of the week.
\item The number of predators had increased by 180\% of the number of predators at the start of the week.
\end{itemize}

The number of prey at the end of the week was \(p\%\) greater than the number of predators at the end of the week.
What is the value of \(p\)?

\vspace{0.75em}

\choicebox{A}{9}
\choicebox{B}{90}
\choicebox{C}{900}
\choicebox{D}{9000}

\vspace{2em}
\newpage

\noindent\markforreview
\vspace{1em}

The exponential function \(f\) is defined by
\[
f(x) = ab^x,
\]
where \(a\) and \(b\) are positive constants.
If \(f(s) = t\) and \(f(s + 1) = t - 0.87t\), where \(s\) and \(t\) are constants, what is the value of \(b\)?

\vspace{0.75em}

\choicebox{A}{0.13}
\choicebox{B}{0.87}
\choicebox{C}{1.13}
\choicebox{D}{1.87}

\vspace{2em}
\newpage

\noindent\markforreview
\vspace{1em}

A beekeeper's initial observation of the population of a certain bee colony was 1,100 bees.
The beekeeper set a goal to increase the population to 2,600 bees.

The beekeeper uses a model that predicts the population of this bee colony:

\begin{itemize}
\item begins at 1,100
\item increases by 120 bees per week in the first two weeks
\item then increases by 180 bees per week until the goal is reached
\end{itemize}

According to this model, at the end of week \(w\) after the initial observation, where \(w > 2\),
which of the following functions gives the predicted number of bees still needed to reach the beekeeper's goal?

\vspace{0.75em}

\choicebox{A}{\(p(w) = 2600 - 180w\)}
\choicebox{B}{\(p(w) = 2480 + 180w\)}
\choicebox{C}{\(p(w) = 1620 - 180w\)}
\choicebox{D}{\(p(w) = -120 + 180w\)}

\vspace{2em}
\newpage

\noindent\markforreview
\vspace{1em}

Given the quadratic equation:
\[
7rx^2 - 28sx + 7r = 0
\]
In the equation:

\begin{itemize}
\item \(r\) is a positive constant
\item \(s\) is a negative constant
\item The equation has two solutions
\end{itemize}

If \(\frac{r}{s} > k\), what is the minimum value of \(k\)?

\vspace{2em}
\newpage

\noindent\markforreview
\vspace{1em}

A square map has a side length of 55 inches, and 1 inch on the map represents an actual distance of 13 miles.
A larger version of the same map is printed as a square with the side length 90\% longer than the side length of the previous map.
On the larger map, which of the following is closest to the actual distance, in miles, represented by 1 inch?

\vspace{0.75em}

\choicebox{A}{1.30}
\choicebox{B}{6.84}
\choicebox{C}{24.70}
\choicebox{D}{49.50}

\vspace{2em}
\newpage

\noindent\markforreview
\vspace{1em}

In isosceles triangle \(ABC\), sides \(AB\) and \(AC\) are congruent.
Point \(D\) divides side \(\overline{BC}\) such that the length of \(\overline{BD}\) is \(\frac{4}{7}\) of the length of \(\overline{BC}\).
Point \(E\) lies on side \(AB\) and point \(F\) lies on side \(AC\) such that when segments \(\overline{DE}\) and \(\overline{DF}\) are drawn, \(\angle BED\) is congruent to \(\angle CFD\).
If the length of \(\overline{BE}\) is 8, what is the length of \(\overline{CF}\)?

\vspace{0.75em}

\choicebox{A}{6}
\choicebox{B}{14}
\choicebox{C}{32}
\choicebox{D}{56}

\vspace{2em}
\newpage

\noindent\markforreview
\vspace{1em}

\[
\frac{\sqrt[3]{x^{17}y^5}}{5x^2\left(\sqrt[8]{y^8}\right)}
\]

For all positive values of \(x\) and \(y\), the given expression is equivalent to which of the following?

\begin{itemize}
\item[I] \(\dfrac{x^{\frac{14}{3}}y^{\frac{5}{3}}}{5xy}\)
\item[II] \(\dfrac{\sqrt[3]{x^{11}y^2}}{5}\)
\end{itemize}

\vspace{0.75em}

\choicebox{A}{I only}
\choicebox{B}{II only}
\choicebox{C}{I and II}
\choicebox{D}{Neither I nor II}

\vspace{2em}
\newpage

\noindent\markforreview
\vspace{1em}

The table shows three values of \(x\) and their corresponding values of \(y\), where \(s\) is a constant.
There is a linear relationship between \(x\) and \(y\).
Which of the following equations represents this relationship?

\[
\begin{array}{|c|c|}
\hline
x & y \\
\hline
-2s & 24 \\
-s & 21 \\
s & 15 \\
\hline
\end{array}
\]

\vspace{0.75em}

\choicebox{A}{\(sx + 3y = 18s\)}
\choicebox{B}{\(3x + sy = 18s\)}
\choicebox{C}{\(3x + sy = 18\)}
\choicebox{D}{\(sx + 3y = 18\)}

\vspace{2em}
\newpage

\noindent\markforreview
\vspace{1em}

The cost of renting a carpet cleaner is \$52 for the first day and \$26 for each additional day.
Which of the following functions gives the cost \(C(d)\), in dollars, of renting the carpet cleaner for \(d\) days, where \(d\) is a positive integer?

\vspace{0.75em}

\choicebox{A}{\(C(d) = 26d + 26\)}
\choicebox{B}{\(C(d) = 26d + 52\)}
\choicebox{C}{\(C(d) = 52d - 26\)}
\choicebox{D}{\(C(d) = 52d + 78\)}

\vspace{2em}
\newpage

\noindent\markforreview
\vspace{1em}

The speed of a vehicle is increasing at a rate of 7.3 meters per second squared.
What is this rate, in miles per minute squared, rounded to the nearest tenth?
(Use 1 mile = 1,609 meters.)

\vspace{0.75em}

\choicebox{A}{0.3}
\choicebox{B}{16.3}
\choicebox{C}{195.8}
\choicebox{D}{220.4}

\vspace{2em}
\newpage

\noindent\markforreview
\vspace{1em}

The table shows three values of \(x\) and their corresponding values of \(g(x)\), where
\[
g(x) = \frac{f(x)}{x + 5}
\]
and \(f\) is a linear function.
What is the \(y\)-intercept of the graph of \(y = f(x)\) in the \(xy\)-plane?

\[
\begin{array}{|c|c|}
\hline
x & g(x) \\
\hline
-20 & 4 \\
-8 & 0 \\
10 & 6 \\
\hline
\end{array}
\]

\vspace{0.75em}

\choicebox{A}{\((0, -8)\)}
\choicebox{B}{\((0, 5)\)}
\choicebox{C}{\((0, 8)\)}
\choicebox{D}{\((0, 40)\)}

\vspace{2em}
